\subsection{\problemblock{*Prop} Data Block}\label{sec:block-prop}

Each \problemblock{*prop} data block defines a thrust force and its corresponding
fuel flow or mass-flow rate. If more than one \problemblock{*prop} block is given
in a segment, the propulsive forces and mass-flow rates from each block are added
together to obtain the total propulsive force and mass-flow rate. If no
\problemblock{*prop} blocks are given, thrust and mass flow are set to zero.

A propulsive-force vector is defined by a magnitude and a direction. The direction is
given by the two angles \statevar{ep1} and \statevar{ep2}, shown in
\cref{fig:thrust-vector-angles} as $\epsilon_1$ and $\epsilon_2$. If both thrust-vector
angles are zero, the default, the thrust force is aligned with the body $x$ axis.

\begin{figure}[htbp]
  \centering
  \begin{tikzpicture}[>=Latex,font=\small,scale=0.9]
  \coordinate (O) at (0,0);
  \draw[->,thick] (O) -- (5.2,0) node[right] {$\hat{x}_b$};
  \draw[->,thick] (O) -- (-0.3,2.0) node[above] {$\hat{z}_b$};
  \draw[->,thick] (O) -- (-1.1,-1.1) node[below left] {$\hat{y}_b$};
  \draw[->,very thick] (O) -- (4.1,1.35) node[above right] {$\vec F_{\mathrm{prop}}$};
  \draw[dashed] (4.1,1.35) -- (4.1,0);
  \draw[->] (1.5,0) arc[start angle=0,end angle=18,radius=1.5];
  \node at (1.7,0.3) {$\epsilon_1$};
  \draw[->] (4.1,0.55) arc[start angle=-90,end angle=-30,radius=0.8];
  \node[right] at (4.65,0.45) {$\epsilon_2$};
  \node[below] at (2.4,-0.35) {Projection of $\vec F_{\mathrm{prop}}$ onto the $\hat{x}_b$--$\hat{z}_b$ plane};
\end{tikzpicture}

  \caption{Thrust-vector angles.}
  \label{fig:thrust-vector-angles}
\end{figure}

Values for the thrust force, thrust-vector angles, and mass-flow rate can be constants
or can be obtained from tables. The following block sets the thrust and mass-flow
rate to constant values:

\begin{lstlisting}[style=taosmanual]
*prop thrust=250 mdot=120.0
      thr_units=kn mdt_units=kg/sec
\end{lstlisting}

The variables \statevar{thrust} and \statevar{mdot} give the thrust force and the
mass-flow rate. The variables \texttt{thr\_units} and \texttt{mdt\_units} give their
units. If no value or table name is provided for a variable, it is set to zero. Thus,
\statevar{ep1} and \statevar{ep2} are both zero in this example.

The default unit for thrust is pounds force, and the default unit for mass-flow rate is
pounds mass per second. Other allowable units are listed in
\cref{tab:propulsion-units}.

\begin{table}[htbp]
  \centering
  \caption{Thrust and mass-flow-rate units.}
  \label{tab:propulsion-units}
  \begin{tabular}{@{}llll@{}}
    \toprule
    Thrust & \multicolumn{3}{c}{Mass flow} \\
    \cmidrule(l){2-4}
    & per second & per minute & per hour \\
    \midrule
    \texttt{lb} & \texttt{lb/sec} & \texttt{lb/min} & \texttt{lb/hr} \\
    \texttt{n}  & \texttt{slugs/sec} & \texttt{slugs/min} & \texttt{slugs/hr} \\
    \texttt{kn} & \texttt{g/sec} & \texttt{g/min} & \texttt{g/hr} \\
                 & \texttt{kg/sec} & \texttt{kg/min} & \texttt{kg/hr} \\
    \bottomrule
  \end{tabular}
\end{table}

Another example,

\begin{lstlisting}[style=taosmanual]
*prop thrust=(thr-recruits) mdot=(mdt-recruits)
      nrecruits=2
\end{lstlisting}

illustrates the use of tables to compute thrust and mass-flow rate. Table
\tableid{thr-recruits} supplies thrust. Table \tableid{mdt-recruits} supplies mass
flow.

The tables require a value for the user-defined variable \texttt{nrecruits}, which
represents the number of motors. Like other data blocks that permit tables to be
specified, a \problemblock{*prop} block may also provide values for user-defined
variables. Once a value has been assigned, it remains set until changed in a later data
block. It need not be repeated in every segment when it does not change.

Units for thrust and mass-flow rate can be specified in the tables, so they are not
shown in this example. If units are given in both places, the units in the
\problemblock{*prop} data block override those in the tables.

Some jet engines have thrust-vector control that can improve takeoff, landing, and
turning performance. For example,

\begin{lstlisting}[style=taosmanual]
*prop thrust=(thr-f100x) mdot=(mdt-f100x)
      ep1=10.0
\end{lstlisting}

sets the first thrust-vector angle to $10^\circ$. In this version of TAOS the thrust-vector
angles can be constants or can be computed with tables; they must be known before
the flight path is computed. The source notes that in future versions they would be
control variables so that they could be varied to meet desired flight conditions.
